% =====================================================================
%  全国大学生数学建模竞赛（CUMCM）官方格式模板 main_template.tex
%  ---------------------------------------------------------------------
%  已按国赛要求修复：
%   1) 删除 \tableofcontents（国赛明确要求论文不得有目录页）；
%   2) 删除作者行 —— 全文任何位置不得出现校名/队号/姓名等身份信息；
%   3) 摘要独立成页：\maketitle 后紧跟摘要，abstract 结束 \newpage；
%   4) 摘要含"关键词"行（分号分隔，5 个左右）；
%   5) 正文主节用"一、二、三、…"中文编号（见下方 ctexset 说明）；
%   6) 页码自摘要页起算（元信息/正文不单独编号）。
%  Verified: compiles with tectonic on 2026-08-23
% =====================================================================

\documentclass[12pt,a4paper]{article}
\usepackage{ctex}
\usepackage{amsmath,amssymb,amsfonts,amsthm}
\usepackage{graphicx}
\usepackage{booktabs}
\usepackage{geometry}
\usepackage{caption}
\usepackage{subcaption}
\usepackage{float}
\usepackage{xcolor}
\usepackage{listings}
\usepackage{longtable}
\usepackage{array}
% hyperref 置于最后，配合可取消目录 (本模板不使用 hyperref 目录，仅用于彩色链接)
\usepackage[hidelinks]{hyperref}

\geometry{left=2.5cm,right=2.5cm,top=2.5cm,bottom=2.5cm}

% ---------------------------------------------------------------------
% 主节编号规范（国赛正文主节用"一、二、三、…"）：
%   方式一（推荐，自动编号）：启用下方两行，则 \section{问题重述} 自动变为"一、问题重述"。
%   方式二（手工编号）：注释掉 ctexset，并在各 \input 分节内用
%       \section*{一、问题重述} 手工写入"一、"等中文编号。
% 任选其一，勿混用。默认开启方式二（各 sections/*.tex 占位文件即采用方式二）。
% ---------------------------------------------------------------------
% \ctexset{section={number=\chinese{section}}}

% 摘要环境排版设置
%\newenvironment{abstract}{%
%    \par\centering\makebox{摘要}\par\nobreak\noindent\rule{\textwidth}{0.4pt}\par\nobreak
%    \small
%}{\par\nobreak\noindent\rule{\textwidth}{0.4pt}\par}

% Listing setting for code appendix
\lstset{
    basicstyle=\footnotesize\ttfamily,
    numbers=left,
    numberstyle=\tiny\color{gray},
    keywordstyle=\color{blue},
    commentstyle=\color{green!60!black},
    stringstyle=\color{orange},
    frame=single,
    breaklines=true,
    backgroundcolor=\color{gray!5}
}

% 主标题占位示例（替换方括号内容即可；国赛标题风格如"基于××模型的××问题研究"）
\title{\textbf{\LARGE 基于×××模型的×××问题研究与决策优化}}
%\date{\today}   % 注意：国赛要求匿名，正文中不宜打印日期/作者；如组委会不禁止可自定义

\begin{document}

% ===================== 摘要页（独立成页） =====================
\maketitle

\begin{abstract}
[此处由各章节摘要汇聚，全面概括各问题的数学解法、核心定理、创新机制与最终结论（约 300 字）。]

\noindent\textbf{关键词}：关键词1；关键词2；关键词3；关键词4；关键词5
\end{abstract}

% 摘要页结束，正文从新页起算
\newpage

% ===================== 正文主节（模块化总装） =====================
% 每个分节文件内使用 \section*{一、……} 手工中文编号；页数预算为建议值。

% 第 1 章  问题背景 / 数据处理与探索性分析
%    应含：问题重述要点、符号说明（表格）、数据可视化（matplotlib/Excel 图）。
%    页数预算：约 2–3 页。
% ===== 第一章 问题重述、符号说明与数据探索（约3页） =====
\section*{一、问题重述与符号说明}

\subsection{问题重述}
“板凳龙”是一种以多节板凳首尾铰接而成的民间舞龙道具。本题中龙头由一节长 341cm 的板凳（含前、后两个把手），龙身与龙尾共 222 节长 220cm 的板凳组成。相邻板凳通过把手销钉铰接，各把手中心沿一条阿基米德螺线分布，螺距记作 $p$。龙头前把手以恒定速率 $1\,\text{m/s}$ 沿螺线向前行进（盘入），其余把手随链运动。总链长（把手中心沿螺线的弧长闭合链）为 $369.16$m。

本题共五问，本文逐一求解：[log:problem]
\begin{enumerate}
\item \textbf{问题一}：螺距 $p=0.55$m，龙头前把手初始位于螺线第 16 圈，求 $t=0\thicksim300$s 内各把手中心的逐秒位置与速度，并列指定节点表；
\item \textbf{问题二}：判定盘入终止（碰撞）时刻；
\item \textbf{问题三}：求使龙头前把手能盘入到以螺线中心为圆心、直径 9m 调头空间的边界（半径 4.5m）的最小螺距；
\item \textbf{问题四}：螺距 1.7m，构造“盘入—S形调头—中心对称盘出”全过程 $-100\thicksim100$s 的位置与速度，并讨论调头圆弧可否缩短；
\item \textbf{问题五}：沿运动路径，求使各把手速度均不超过 2\,m/s 的龙头最大行进速度。
\end{enumerate}

\subsection{符号说明}
本文所用主要符号见表 \ref{tab:sym}。[log:problem, scale]

\begin{table}[H]\centering
\caption{主要符号说明[log:scale]}
\label{tab:sym}\begin{tabular}{lll}
\toprule
符号 & 意义 & 值/单位 \\
\midrule
$p$ & 螺距 & 0.55m(Q1/2)、1.7m(Q4)、0.30m(Q3) \\
$a$ & 螺线尺度 $a=p/2\pi$ & m/rad \\
$r,\theta$ & 螺线极坐标 & m、rad \\
$L_h,L_b$ & 龙头/龙身孔距 & 2.86m、1.65m \\
$w$ & 板凳板宽 & 0.30m \\
$N$ & 把手中心数 & 224 \\
$D_i$ & 把手 $i$ 距龙头前把手的弧长偏移 & 0,…,369.16m \\
$s(\theta)$ & 螺线弧长闭式 & m \\
$S_0$ & 龙头初始弧长 $s(32\pi)$ & 442.59026m \\
$t^{\ast}$ & Q2 终止时刻 & 73.430s \\
$p^{\ast}$ & Q3 最小螺距 & 0.30m \\
$R_1,R_2$ & Q4 S形圆弧半径 & 3.0054m、1.5027m \\
$\beta_{\max}$ & Q5 龙头最大速度 & 2.0\,m/s \\
\bottomrule
\end{tabular}
\end{table}

\subsection{数据探索与自变量确定}
本题无外部数据附件，全部参数由题面几何推导（板长 $341$cm、$220$cm，孔距等于板长减去 $2\times27.5$cm 的孔径偏置，故龙头孔距 $2.86$m、龙身孔距 $1.65$m）。初始龙头前把手位于“第16圈”外端，本文取 $\theta_0=32\pi$（$r_0=8.80$m），该选定在问题一假设中声明。$\theta_0$ 的敏感性在第七章灵敏度分析中给出。独立复核显示，龙头 $r=8.80$m 时，等弧距铺排下龙尾后把手位于 $r\approx3.58$m，初始盘面覆盖螺线第 6.5--16 圈，无自相交。[log:q1.s0]

% 第 2 章  问题（一）模型建立与求解
%    应含：模型假设与符号、数学模型公式、求解算法、结果表、结果分析。
%    页数预算：约 3–4 页。
% ===== 第 2 章实际内容：问题（一）模型建立与求解（约 3–4 页） =====
% 在此写入：模型假设与符号、数学模型公式、求解算法、结果表、结果分析。
\section*{二、问题（一）模型的建立与求解}
\noindent（占位章节：在此补充第二问的完整建模内容。）

% 第 3 章  问题（二）模型建立与求解（如含风险/不确定性建模、CVaR 等）
%    应含：扩展模型、核心定理或推导、数值求解、结果对比。
%    页数预算：约 3–4 页。
% ===== 第 3 章实际内容：问题（二）模型（含风险/不确定性、CVaR 等，约 3–4 页） =====
% 在此写入：扩展模型、核心定理或推导、数值求解、结果对比。
\section*{三、问题（二）模型的建立与求解}
\noindent（占位章节：在此补充第三问的完整建模内容。）

% 第 4 章  问题（三）模型建立与求解（如相关性、多目标/非线性扩展）
%    应含：进一步建模、指标体系、求解与结果图/表。
%    页数预算：约 3–4 页。
% ===== 第 4 章实际内容：问题（三）模型（相关性/多目标/非线性扩展，约 3–4 页） =====
% 在此写入：进一步建模、指标体系、求解与结果图/表。
\section*{四、问题（三）模型的建立与求解}
\noindent（占位章节：在此补充第四问的完整建模内容。）

% 第 5 章  模型评价、灵敏度分析、改进方向与附录
%    应含：灵敏度/鲁棒性分析、优缺点与改进、模型推广；代码附录用 listings。
%    页数预算：约 2–3 页。
% ===== 第 5 章实际内容：模型评价、灵敏度分析、改进与附录（约 2–3 页） =====
% 在此写入：灵敏度/鲁棒性分析、优缺点与改进、模型推广；代码附录用 listings。
\section*{五、模型的评价、改进与附录}
\noindent（占位章节：在此补充灵敏度分析与模型评价内容。）

\end{document}